\chapter{Developmental and Behavioral Screening Tests}
\section*{What Is This Test?} Developmental and behavioral screening checks whether a child is meeting expected movement, language, learning, social, and emotional milestones.
\section*{Alternative Names} Developmental surveillance; milestone screening; behavioral assessment.
\section*{Principle} Validated questionnaires and observation compare reported skills with age- and development-appropriate expectations.
\section*{Why This Test Is Ordered} Screening is routine at selected well-child visits and whenever parents or clinicians notice delay, regression, behavior concerns, or difficulty functioning.
\section*{Primary vs Secondary Test} Screening identifies concern; it does not diagnose a developmental disorder. Diagnostic evaluation is the secondary next step.
\section*{How This Test Is Performed} Caregivers complete questions and the clinician observes the child, reviews history, and discusses behavior across settings.
\section*{What Preparation Is Needed} Bring developmental, school, therapy, hearing, vision, and medical information.
\section*{Understanding Results} A positive result means further assessment is needed; a negative screen does not override persistent concern or regression.
\section*{Follow-up Tests} Hearing and vision testing, speech-language, cognitive, autism, occupational, or specialist evaluation may follow.
\section*{References} \begin{enumerate}\item MedlinePlus. Developmental and Behavioral Screening Tests.\item American Academy of Pediatrics. Developmental surveillance and screening guidance.\end{enumerate}
\clearpage\section*{Understanding Results and Follow-up}
The laboratory report should identify the specimen, method, units, reference interval or decision limit, and whether the result is positive, negative, abnormal, or inconclusive. Reference intervals vary with age, sex, pregnancy, specimen type, assay, and laboratory; a result outside a range is not automatically a diagnosis, and a result within a range does not always exclude disease. Discuss unexpected results with the ordering clinician, who can decide whether repeat or confirmatory testing, treatment, or referral is appropriate. Do not change medicines or delay urgent care based on an isolated result.
\section*{Regulatory Status and Authorization Date}This chapter describes a test category, not a specific commercial product. FDA, CDSCO, EU IVDR, and MHRA approval or registration dates therefore cannot be assigned without the assay manufacturer, product name, instrument, and intended use. For laboratory-developed tests, an individual agency approval date may not exist. See the Preface for the interpretation of regulatory status.
\clearpage
